\rok{Фебруар 2026.}{Ф}

Други практични тест из Алгоритама и структура податак

Први практични тест одржан 13.02.2026.

\zadatak{1}{Имплементација DFSVisit методе Depth-First-Search алгоритма}
Написати имплементацију DFSVisit методе истоимене алгоритма.

\textbf{Додатне напомене за решавање задатка:}
\begin{itemize}
	\item a) Написати алгоритам.
	\item b) У Main методи креирати произволан граф. Обезбедити начин за тестирање и проверу нове функционалности, односно позив DFS методе.
	\item У template-у је метода \texttt{private void DFSVisit(LinkedList.Node node, int nodeIndex)} која је основна метода за рекурзивно извршавање DFS алгоритма.
\end{itemize}

\zadatak{2}{Карактери са тачним бројем појављивања}
Написати методу која прима string str и цео број k. Метода треба да исписе все карактере из string-а који се у њему појављују тачно k пута, користећи хеш табелу ради ефикасности.

\textbf{Додатни захтеви за решавање задатка:}
\begin{itemize}
	\item Користити Hash табелу за чување броја појављивања (као кључ је карактер, као вредност је број).
	\item Имплементацију писати у методи:
	\begin{Verbatim}[fontsize=\small]
public static void printCharsWithCountK(string str, int k)
	\end{Verbatim}
	\item Разлика између малих и великих слова постоји.
\end{itemize}

\textbf{Примери:}
\begin{itemize}
	\item \textbf{Улаз 1:} \texttt{str = "Algoritmi i strukture", k = 3}
	\item \textbf{Излаз 1:} \texttt{"i"}
	\item \textbf{Улаз 2:} \texttt{str = "Vetar", k = 2}
	\item \textbf{Излаз 2:} \texttt{"Nijedan karakter se ne pojavljuje 2 puta"}
\end{itemize}

\zadatak{3}{Сума листова већа од X}
Написати методу која проверава да ли је сума свих листова у стаблу већа од заданте вредности threshold.

\textbf{Додатни захтеви за решавање задатка:}
\begin{itemize}
	\item Јавна метода:
	\begin{Verbatim}[fontsize=\small]
public bool isLeafSumGreater(int threshold)
	\end{Verbatim}
	\item Приватна метода:
	\begin{Verbatim}[fontsize=\small]
private int getLeafSum(Node current)
	\end{Verbatim}
	\item Подразумева се да стабло има барем два чвора.
	\item Алгоритам ОБАВЕЗНО писати у оквиру класе \texttt{BinarySearchTree}.
\end{itemize}

\textbf{Примери:}

\begin{center}
\begin{minipage}{\textwidth}
\centering
\begin{forest}
for tree={
	circle,
	draw,
	thick,
	fill=gray!20,
	minimum size=8mm,
	inner sep=1pt,
	s sep=8mm,
	l sep=12mm,
	edge={-, thick},
	font=\small
}
[8
	[4
		[2
			[, phantom]
			[3]
		]
		[5]
	]
	[9
		[, phantom]
		[11]
	]
]
\end{forest}

\vspace{0.3cm}
\textbf{Слика 1.} Пример BST-a.
\end{minipage}
\end{center}

\begin{itemize}
	\item \textbf{Улаз 1:} \texttt{threshold = 20}
	\item \textbf{Излаз 1:} \texttt{false (suma listova = 19 (3 + 5 + 11))}
\end{itemize}

\begin{center}
\begin{minipage}{\textwidth}
\centering
\begin{forest}
for tree={
	circle,
	draw,
	thick,
	fill=gray!20,
	minimum size=8mm,
	inner sep=1pt,
	s sep=8mm,
	l sep=12mm,
	edge={-, thick},
	font=\small
}
[10
	[4
		[2
			[3]
			[, phantom]
		]
		[6
			[5]
			[9]
		]
	]
	[20
		[12
			[, phantom]
			[14]
		]
		[29
			[27]
			[, phantom]
		]
	]
]
\end{forest}

\vspace{0.3cm}
\textbf{Слика 2.} Пример BST-a.
\end{minipage}
\end{center}

\begin{itemize}
	\item \textbf{Улаз 2:} \texttt{threshold = 20}
	\item \textbf{Излаз 2:} \texttt{true (suma listova = 58 (3 + 5 + 9 + 14 + 27))}
\end{itemize}

\sablon{Шаблон другог теста фебруар 2026. група Ф}{2025/Februar/GrupaF/AISP2025_Test2_GrupaF.linq}
